%% toc: P2 22 22.1,22.2
\section{Proofs for the one-sorted reduction at every cardinality}
\label{sec:annex-p02-w02b-p02-033-any-cardinality}

Throughout, \(\Lang_{\mathrm{MS}}\) is a strict many-sorted signature whose
sets \(\mathsf{Sort}\), \(\mathcal F\) and \(\mathcal R\) may have any
cardinality, with the variable conventions of
\cref{def:p02-w02b-p02-033-variables}.  For a sorted structure
\(\mathcal M\), a total assignment \(g\), a sorted variable \(v\) and
\(b\in M_{\operatorname{srt}(v)}\), \(g[v\mapsto b]\) is as in
\cref{def:p02-w02-p02-038-sorted-structure}; for a one-sorted assignment
\(\sigma\), \(\sigma[x\mapsto d]\) sends \(x\) to \(d\) and agrees with
\(\sigma\) elsewhere.  One-sorted satisfaction uses the clauses of
\cref{sec:fol-semantics-02}: \(\mathcal B,\sigma\entails\forall x\,\chi\) iff
\(\mathcal B,\sigma[x\mapsto d]\entails\chi\) for every \(d\in\dom\mathcal B\),
dually for \(\exists\), and the Boolean clauses are the truth tables.  The
metatheory is ZFC, and no proof below uses the Axiom of Choice.

\begin{annexitem}{annex:p02-w02b-p02-033-sentences}{thm:p02-w02b-p02-033-sentences}{Parameter names are fresh, and translations are sentences}
\begin{proof}
\emph{Clause 1.}  If \(\mathsf n_v=\mathsf n_w\), then \((W,v)=(W,w)\), and
equality of ordered pairs gives \(v=w\).  Suppose
\(\mathsf n_v=(W,v)=\{\{W\},\{W,v\}\}\in W\).  Then
\(W\in\{W\}\in(W,v)\in W\).  The set \(Z\defeq\{W,\{W\},(W,v)\}\) has no
\(\in\)-minimal element: \((W,v)\in W\cap Z\), \(W\in\{W\}\cap Z\) and
\(\{W\}\in(W,v)\cap Z\).  This contradicts Foundation, so
\(\mathsf n_v\notin W\).  Hence the names are pairwise distinct and disjoint
from \(\mathcal F\), \(\mathcal R\), the predicates \(Q_{\mathsf s}\),
\(\mathsf{Var}\) and the logical symbols, which is the disjointness required
by \cref{def:p02-fol-signature}.

\emph{Clause 2, terms.}  By induction on \(t\).  If \(t\) is a variable
\(v\in\dom\eta\), then \(t^{\eta}=\eta(v)\in\operatorname{ran}\eta\).  If
\(t\) is a variable \(v\notin\dom\eta\), then \(t^{\eta}=\mathsf n_v\) has no
variable.  A constant has no variable.  The variables of
\(f(t_1,\ldots,t_n)^{\eta}=f(t_1^{\eta},\ldots,t_n^{\eta})\) are those of the
\(t_i^{\eta}\), which lie in \(\operatorname{ran}\eta\) by the induction
hypothesis.

\emph{Clause 2, formulas.}  By induction on \(\varphi\), simultaneously for
all environments.  The free variables of a translated atom are the variables
of its translated terms, which lie in \(\operatorname{ran}\eta\) by the term
case; \(\bottom\) has none; and the free variables of
\(\neg\chi\) and of \(\chi\mathbin{\square}\chi'\) are those of the
components, so the induction hypothesis at the same \(\eta\) applies.  For
\(\forall v\,\varphi\) put \(k\defeq k(\eta)\) and
\(\eta'\defeq\eta\langle v\rangle\).  The free variables of
\(\forall x_k\bigl(Q_{\operatorname{srt}(v)}(x_k)\to
\operatorname{tr}^{\mathsf n}_{\eta'}(\varphi)\bigr)\) are those of
\(\operatorname{tr}^{\mathsf n}_{\eta'}(\varphi)\) other than \(x_k\), since
\(Q_{\operatorname{srt}(v)}(x_k)\) contributes only \(x_k\).  By the induction
hypothesis they lie in \(\operatorname{ran}\eta'\setminus\{x_k\}\).  Every
value of \(\eta'\) other than \(x_k\) is a value \(\eta(u)\) with
\(u\in\dom\eta\setminus\{v\}\), so
\(\operatorname{ran}\eta'\setminus\{x_k\}\subseteq\operatorname{ran}\eta\).
The case of \(\exists v\,\varphi\) is the same computation with \(\land\) in
place of \(\to\) and \(\exists x_k\) in place of \(\forall x_k\), and gives
the same set of free variables.  For \(\eta=\varnothing\) the range is empty,
so \(\operatorname{tr}^{\mathsf n}(\varphi)\) has no free variable.  Its
nonlogical symbols are those of \(\varphi\), the predicates
\(Q_{\mathsf s}\) and names \(\mathsf n_v\), all in
\(\Lang_{\mathrm{MS}}^{\mathsf n}\); so it is an
\(\Lang_{\mathrm{MS}}^{\mathsf n}\)-sentence.
\end{proof}
\end{annexitem}

\begin{annexitem}{annex:p02-w02b-p02-033-isomorphism}{thm:p02-w02b-p02-033-isomorphism}{Sorted isomorphisms preserve satisfaction}
\begin{proof}
\emph{Terms}, by induction on \(t\), for every total \(g\).  For a variable
\(v\) of sort \(\mathsf s\),
\(\llbracket v\rrbracket^{\mathcal N}_{\pi g}=(\pi g)(v)=
\pi_{\mathsf s}(g(v))=\pi_{\mathsf s}(\llbracket v\rrbracket^{\mathcal M}_{g})\).
For a constant \(c\) of sort \(\mathsf s\),
\(\llbracket c\rrbracket^{\mathcal N}_{\pi g}=c^{\mathcal N}=
\pi_{\mathsf s}(c^{\mathcal M})\).  For
\(f:\mathsf s_1\times\cdots\times\mathsf s_n\longrightarrow\mathsf s\) and
\(b_i\defeq\llbracket t_i\rrbracket^{\mathcal M}_{g}\), the induction
hypothesis gives
\(\llbracket t_i\rrbracket^{\mathcal N}_{\pi g}=\pi_{\mathsf s_i}(b_i)\), so
\(\llbracket f(t_1,\ldots,t_n)\rrbracket^{\mathcal N}_{\pi g}=
f^{\mathcal N}(\pi_{\mathsf s_1}(b_1),\ldots,\pi_{\mathsf s_n}(b_n))=
\pi_{\mathsf s}(f^{\mathcal M}(b_1,\ldots,b_n))\).

\emph{Formulas}, by induction on \(\varphi\), for every total \(g\).  For
\(P(t_1,\ldots,t_n)\) with \(b_i\) as above,
\(\mathcal N\entails P(t_1,\ldots,t_n)[\pi g]\) iff
\((\pi_{\mathsf s_1}(b_1),\ldots,\pi_{\mathsf s_n}(b_n))\in P^{\mathcal N}\)
iff \((b_1,\ldots,b_n)\in P^{\mathcal M}\) iff
\(\mathcal M\entails P(t_1,\ldots,t_n)[g]\).  For \(t=u\) with both terms of
sort \(\mathsf s\), the term case gives values \(\pi_{\mathsf s}(a)\) and
\(\pi_{\mathsf s}(b)\) in \(\mathcal N\), where \(a,b\) are the values in
\(\mathcal M\); as \(\pi_{\mathsf s}\) is injective,
\(\pi_{\mathsf s}(a)=\pi_{\mathsf s}(b)\) iff \(a=b\).  The formula
\(\bottom\) is false on both sides.  For \(\neg\chi\), \(\chi\land\chi'\),
\(\chi\lor\chi'\) and \(\chi\to\chi'\), the clause on each side is the same
truth function of the components, which agree by the induction hypothesis at
\(g\).

For \(\forall v\,\chi\) with \(v\) of sort \(\mathsf s\): for every
\(b\in M_{\mathsf s}\), \(\pi(g[v\mapsto b])=(\pi g)[v\mapsto\pi_{\mathsf s}(b)]\),
since both send \(v\) to \(\pi_{\mathsf s}(b)\) and every other \(u\) to
\(\pi_{\operatorname{srt}(u)}(g(u))\).  Hence, by the induction hypothesis at
\(g[v\mapsto b]\),
\(\mathcal M\entails\chi[g[v\mapsto b]]\) iff
\(\mathcal N\entails\chi[(\pi g)[v\mapsto\pi_{\mathsf s}(b)]]\).  Now
\(\mathcal M\entails\forall v\,\chi[g]\) says the left side holds for all
\(b\in M_{\mathsf s}\), and \(\mathcal N\entails\forall v\,\chi[\pi g]\) says
\(\mathcal N\entails\chi[(\pi g)[v\mapsto d]]\) for all \(d\in N_{\mathsf s}\).
These agree because \(b\mapsto\pi_{\mathsf s}(b)\) maps \(M_{\mathsf s}\)
onto \(N_{\mathsf s}\): every \(d\in N_{\mathsf s}\) is
\(\pi_{\mathsf s}(\pi_{\mathsf s}^{-1}(d))\), and every
\(\pi_{\mathsf s}(b)\) lies in \(N_{\mathsf s}\).  For \(\exists v\,\chi\),
\(\mathcal M\entails\exists v\,\chi[g]\) says the left side holds for some
\(b\in M_{\mathsf s}\); then \(d\defeq\pi_{\mathsf s}(b)\in N_{\mathsf s}\)
witnesses \(\mathcal N\entails\exists v\,\chi[\pi g]\).  Conversely a
witness \(d\in N_{\mathsf s}\) on the right gives the witness
\(b\defeq\pi_{\mathsf s}^{-1}(d)\) on the left, because
\(\pi_{\mathsf s}(b)=d\).
\end{proof}
\end{annexitem}

\begin{annexitem}{annex:p02-w02b-p02-033-restriction}{thm:p02-w02b-p02-033-restriction}{Restriction correspondence for the parameter form}
\begin{proof}
Write \(M_{\mathsf s}\defeq Q_{\mathsf s}^{\mathcal B}\) for the carriers of
\(\mathcal B^{\downarrow}\).

\emph{Clause 1.}  For each sort \(\mathsf s\), the parameter guard contains
\(Q_{\mathsf s}(\mathsf n_{x^{\mathsf s}})\), so
\(\mathsf n_{x^{\mathsf s}}^{\mathcal B}\in M_{\mathsf s}\) and
\(M_{\mathsf s}\ne\varnothing\).  \(\mathcal B^{\downarrow}\) is a sorted
structure: \(c^{\mathcal B}\in M_{\mathsf s}\) for \(c:\longrightarrow
\mathsf s\) by \((\mathrm{Ct})\); for
\(f:\mathsf s_1\times\cdots\times\mathsf s_n\longrightarrow\mathsf s\) with
\(n\geq1\) and \(b_i\in M_{\mathsf s_i}\), \((\mathrm{Cl})\) gives
\(f^{\mathcal B}(b_1,\ldots,b_n)\in M_{\mathsf s}\), so the restriction of
\(f^{\mathcal B}\) to \(M_{\mathsf s_1}\times\cdots\times M_{\mathsf s_n}\)
maps into \(M_{\mathsf s}\); and \(P^{\mathcal B^{\downarrow}}\) is a subset
of the declared product by its definition.  If
\(\mathcal B\entails(\mathrm{Dj})\), then for \(\mathsf s\ne\mathsf t\) no
element satisfies both \(Q_{\mathsf s}\) and \(Q_{\mathsf t}\), so the
nonempty carriers are pairwise disjoint and \(\mathcal B^{\downarrow}\) is
strict.  For \(v\in\dom\eta\),
\(\theta(v)=\sigma(\eta(v))\in M_{\operatorname{srt}(v)}\) because \(\sigma\)
respects \(\eta\); for \(v\notin\dom\eta\),
\(\theta(v)=\mathsf n_v^{\mathcal B}\in M_{\operatorname{srt}(v)}\) by the
parameter guard.  So \(\theta\) is a total assignment into
\(\mathcal B^{\downarrow}\).

\emph{Clause 2}, by induction on \(t\).  For a variable \(v\in\dom\eta\),
\(\llbracket\eta(v)\rrbracket^{\mathcal B}_{\sigma}=\sigma(\eta(v))=
\theta(v)=\llbracket v\rrbracket^{\mathcal B^{\downarrow}}_{\theta}\).  For
\(v\notin\dom\eta\),
\(\llbracket\mathsf n_v\rrbracket^{\mathcal B}_{\sigma}=
\mathsf n_v^{\mathcal B}=\theta(v)\).  In both cases the value lies in
\(M_{\operatorname{srt}(v)}\) by clause~1.  For a constant \(c\) of sort
\(\mathsf s\), both values are \(c^{\mathcal B}\in M_{\mathsf s}\).  For
\(f(t_1,\ldots,t_n)\) with \(f\) as above, the induction hypothesis gives
common values \(b_i\in M_{\mathsf s_i}\) of \(t_i^{\eta}\) in \(\mathcal B\)
and of \(t_i\) in \(\mathcal B^{\downarrow}\); since the tuple lies in the
product, \(f^{\mathcal B}(b_1,\ldots,b_n)=
f^{\mathcal B^{\downarrow}}(b_1,\ldots,b_n)\), which lies in
\(M_{\mathsf s}\) by clause~1.

\emph{Clause 3}, by induction on \(\varphi\), simultaneously for every
environment \(\eta\) and every assignment \(\sigma\) respecting it.

\emph{Relation atoms.}  For \(P:\mathsf s_1\times\cdots\times\mathsf s_n\)
let \(b_i\) be the common values from clause~2; each \(b_i\in
M_{\mathsf s_i}\).  Then \(\mathcal B,\sigma\entails
P(t_1^{\eta},\ldots,t_n^{\eta})\) iff \((b_1,\ldots,b_n)\in P^{\mathcal B}\)
iff \((b_1,\ldots,b_n)\in P^{\mathcal B}\cap(M_{\mathsf s_1}\times\cdots
\times M_{\mathsf s_n})=P^{\mathcal B^{\downarrow}}\) iff
\(\mathcal B^{\downarrow}\entails P(t_1,\ldots,t_n)[\theta]\).

\emph{Equations.}  By clause~2, \(t^{\eta}\) and \(u^{\eta}\) have in
\(\mathcal B\) the values that \(t\) and \(u\) have in
\(\mathcal B^{\downarrow}\) at \(\theta\); both sides compare these two
values for identity.

\emph{\(\bottom\) and connectives.}  \(\bottom\) is false on both sides.
For \(\neg\chi\) and \(\chi\mathbin{\square}\chi'\), the translation commutes
with the connective, each side applies the same truth function, and the
components agree by the induction hypothesis at the same \(\eta\) and
\(\sigma\).

\emph{Universal quantifier.}  Let \(\varphi=\forall v\,\chi\),
\(\mathsf s\defeq\operatorname{srt}(v)\), \(k\defeq k(\eta)\) and
\(\eta'\defeq\eta\langle v\rangle\).  By the one-sorted clauses for
\(\forall\) and \(\to\),
\(\mathcal B,\sigma\entails\operatorname{tr}^{\mathsf n}_{\eta}(\varphi)\)
iff \(\mathcal B,\sigma_d\entails\operatorname{tr}^{\mathsf n}_{\eta'}(\chi)\)
for every \(d\in M_{\mathsf s}\), where \(\sigma_d\defeq\sigma[x_k\mapsto d]\).
Fix \(d\in M_{\mathsf s}\).  First, \(\sigma_d\) respects \(\eta'\):
\(\sigma_d(\eta'(v))=\sigma_d(x_k)=d\in M_{\mathsf s}\); and for
\(u\in\dom\eta'\setminus\{v\}\) we have \(u\in\dom\eta\) and
\(\eta'(u)=\eta(u)\ne x_k\), because \(x_k\notin\operatorname{ran}\eta\), so
\(\sigma_d(\eta'(u))=\sigma(\eta(u))\in M_{\operatorname{srt}(u)}\).  Second,
\(\theta_{\eta',\sigma_d}=\theta[v\mapsto d]\): at \(v\) both are \(d\); at
\(u\ne v\) with \(u\in\dom\eta\), both are \(\sigma(\eta(u))\) by the
computation just made; at \(u\ne v\) with \(u\notin\dom\eta\), \(u\) lies
outside \(\dom\eta'\) and both are \(\mathsf n_u^{\mathcal B}\).  By the
induction hypothesis at \(\eta'\) and \(\sigma_d\),
\(\mathcal B,\sigma_d\entails\operatorname{tr}^{\mathsf n}_{\eta'}(\chi)\) iff
\(\mathcal B^{\downarrow}\entails\chi[\theta[v\mapsto d]]\).  So the left side
holds iff \(\mathcal B^{\downarrow}\entails\chi[\theta[v\mapsto d]]\) for
every \(d\in M_{\mathsf s}\), which is the clause for
\(\mathcal B^{\downarrow}\entails\forall v\,\chi[\theta]\).

\emph{Existential quantifier.}  For \(\varphi=\exists v\,\chi\), with the same
notation, the one-sorted clauses for \(\exists\) and \(\land\) give
\(\mathcal B,\sigma\entails\operatorname{tr}^{\mathsf n}_{\eta}(\varphi)\) iff
\(\mathcal B,\sigma_d\entails\operatorname{tr}^{\mathsf n}_{\eta'}(\chi)\) for
some \(d\in M_{\mathsf s}\).  For each such \(d\) the two facts established
in the universal case (\(\sigma_d\) respects \(\eta'\), and
\(\theta_{\eta',\sigma_d}=\theta[v\mapsto d]\)) did not depend on the
quantifier, so the induction hypothesis gives the equivalence with
\(\mathcal B^{\downarrow}\entails\chi[\theta[v\mapsto d]]\).  So the left side
holds iff some \(d\in M_{\mathsf s}\) satisfies
\(\mathcal B^{\downarrow}\entails\chi[\theta[v\mapsto d]]\), which is the
clause for \(\mathcal B^{\downarrow}\entails\exists v\,\chi[\theta]\).
\end{proof}
\end{annexitem}

\begin{annexitem}{annex:p02-w02b-p02-033-tagging}{thm:p02-w02b-p02-033-tagging}{Tagging correspondence for the parameter form}
\begin{proof}
\emph{The structures are well defined.}  If \(\star\in D_0\), then
\((\mathsf{Sort},\mathsf{Sort})=(\mathsf s,a)\) for some
\(\mathsf s\in\mathsf{Sort}\), so \(\mathsf{Sort}=\mathsf s\in\mathsf{Sort}\),
and \(\{\mathsf{Sort}\}\) has no \(\in\)-minimal element, contradicting
Foundation.  So \(\star\notin D_0\).  \(D_0\) is nonempty because
\(\mathsf{Sort}\) and every carrier are nonempty.  An element
\((\mathsf s_i,b_i)\) determines \(b_i\), so the first case of each function
is single-valued; the other values, \(\star\) or
\((\mathsf s,g(x^{\mathsf s}))\), lie in the respective domains.  Each
\(Q_{\mathsf s}\), constant, name and relation is interpreted inside the
domain.

\emph{Clause 1.}  \((\mathrm{Ne})\): \((\mathsf s,g(x^{\mathsf s}))\in
Q_{\mathsf s}^{\mathcal B}\).  \((\mathrm{Ct})\): \(c^{\mathcal B}=
(\mathsf s,c^{\mathcal M})\in\{\mathsf s\}\times M_{\mathsf s}\).
\((\mathrm{Cl})\): if \(d_i\in Q_{\mathsf s_i}^{\mathcal B}\), then
\(d_i=(\mathsf s_i,b_i)\) with \(b_i\in M_{\mathsf s_i}\), the first case
applies, and \(f^{\mathcal B}(d_1,\ldots,d_n)=
(\mathsf s,f^{\mathcal M}(b_1,\ldots,b_n))\in Q_{\mathsf s}^{\mathcal B}\).
\((\mathrm{Dj})\): \((\mathsf s,a)=(\mathsf t,b)\) implies
\(\mathsf s=\mathsf t\), so \(Q_{\mathsf s}^{\mathcal B}\cap
Q_{\mathsf t}^{\mathcal B}=\varnothing\) for \(\mathsf s\ne\mathsf t\).
Parameter guard: \(\mathsf n_v^{\mathcal B}=(\operatorname{srt}(v),g(v))\in
Q_{\operatorname{srt}(v)}^{\mathcal B}\).  Every member of every
\(Q_{\mathsf s}^{\mathcal B}\) lies in \(D_0\), and \(\star\notin D_0\), so
\(\star\) lies in no \(Q_{\mathsf s}\).  If
\(\mathsf{Sort}=\{\mathsf s_1,\ldots,\mathsf s_k\}\), every element
\((\mathsf s,a)\) of \(\mathcal M^{\circ}_{g}\) has
\(\mathsf s=\mathsf s_j\) for some \(j\) and lies in
\(Q_{\mathsf s_j}\), so \(\mathsf{Cov}\) holds; in \(\mathcal M^{\star}_{g}\)
the element \(\star\) falsifies the disjunction, so \(\mathsf{Cov}\) fails.

\emph{Clause 2.}  By clause~1, \(\mathcal B\) satisfies \((\mathrm{Ct})\),
\((\mathrm{Cl})\), \((\mathrm{Dj})\) and the parameter guard, so
\cref{thm:p02-w02b-p02-033-restriction} applies: \(\mathcal B^{\downarrow}\)
is strict with carriers \(\{\mathsf s\}\times M_{\mathsf s}\),
\(\theta\defeq\theta_{\eta,\sigma}\) is a total assignment into it, and
\(\mathcal B,\sigma\entails\operatorname{tr}^{\mathsf n}_{\eta}(\varphi)\) iff
\(\mathcal B^{\downarrow}\entails\varphi[\theta]\).  As
\(\theta(v)\in\{\operatorname{srt}(v)\}\times M_{\operatorname{srt}(v)}\),
it has the stated form with a unique \(a_v\), and \(g_{\eta,\sigma}\) is a
total assignment into \(\mathcal M\).  Let
\(\pi_{\mathsf s}(a)\defeq(\mathsf s,a)\), a bijection of
\(M_{\mathsf s}\) onto \(\{\mathsf s\}\times M_{\mathsf s}\).  It satisfies
the hypotheses of \cref{thm:p02-w02b-p02-033-isomorphism} from \(\mathcal M\)
to \(\mathcal B^{\downarrow}\): \(c^{\mathcal B^{\downarrow}}=
(\mathsf s,c^{\mathcal M})=\pi_{\mathsf s}(c^{\mathcal M})\); for
\(b_i\in M_{\mathsf s_i}\),
\(f^{\mathcal B^{\downarrow}}(\pi_{\mathsf s_1}(b_1),\ldots,
\pi_{\mathsf s_n}(b_n))=(\mathsf s,f^{\mathcal M}(b_1,\ldots,b_n))=
\pi_{\mathsf s}(f^{\mathcal M}(b_1,\ldots,b_n))\); and
\(P^{\mathcal B^{\downarrow}}=P^{\mathcal B}\), since \(P^{\mathcal B}\) already
lies in the product of the carriers, while
\((\pi_{\mathsf s_1}(b_1),\ldots,\pi_{\mathsf s_n}(b_n))\in P^{\mathcal B}\) iff
it equals a tuple \(((\mathsf s_1,b'_1),\ldots,(\mathsf s_n,b'_n))\) with
\((b'_1,\ldots,b'_n)\in P^{\mathcal M}\), iff
\((b_1,\ldots,b_n)\in P^{\mathcal M}\), because the tagged pairs determine
\(b'_i=b_i\).  Since \(\pi g_{\eta,\sigma}=\theta\) by the definition of
\(a_v\), the lemma gives \(\mathcal M\entails\varphi[g_{\eta,\sigma}]\) iff
\(\mathcal B^{\downarrow}\entails\varphi[\theta]\).  Combining the two
equivalences proves clause~2.

\emph{Clause 3.}  Every assignment respects \(\varnothing\).  For
\(\eta=\varnothing\), \(\theta_{\varnothing,\sigma}(v)=\mathsf n_v^{\mathcal B}
=(\operatorname{srt}(v),g(v))\), so \(g_{\varnothing,\sigma}=g\).  Clause~2
gives \(\mathcal B,\sigma\entails\operatorname{tr}^{\mathsf n}(\varphi)\) iff
\(\mathcal M\entails\varphi[g]\), and for a strict structure and total
assignment the latter is \(\mathcal M,g\entails\varphi\)
(\cref{def:p02-w02-p02-038-satisfaction}).
\end{proof}
\end{annexitem}

\begin{annexitem}{annex:p02-w02b-p02-033-reduction}{thm:p02-w02b-p02-033-reduction}{Reduction of strict many-sorted consequence for every strict signature}
\begin{proof}
Write \(\mathsf{Ax}^{-}\defeq(\mathrm{Ct})\cup(\mathrm{Cl})\cup(\mathrm{Dj})\)
for the sort theory without \((\mathrm{Ne})\), and let \(X\) be any of the
four premise sets of clauses 1--3: \(\mathsf{Ax}_{\mathrm{sort}}\) or
\(\mathsf{Ax}^{-}\), without \(\mathsf{Cov}\) or, when \(\mathsf{Sort}\) is
finite, with it, together with \(\mathsf{ParGuard}\) and
\(\operatorname{tr}^{\mathsf n}[\Gamma]\).  In each case
\(\mathsf{Ax}^{-}\cup\mathsf{ParGuard}\cup
\operatorname{tr}^{\mathsf n}[\Gamma]\subseteq X\), and \(X\) is contained in
\(\mathsf{Ax}_{\mathrm{sort}}\cup\mathsf{ParGuard}\cup
\operatorname{tr}^{\mathsf n}[\Gamma]\), with \(\mathsf{Cov}\) added when
\(X\) contains it.  It suffices to prove, for each such \(X\), that
\(\Gamma\entails_{\mathrm{MS}}\varphi\) iff
\(X\entails_{\mathrm{FOL}}\operatorname{tr}^{\mathsf n}(\varphi)\).

\emph{From many-sorted to one-sorted.}  Suppose
\(\Gamma\entails_{\mathrm{MS}}\varphi\), and let \(\mathcal B,\sigma\) satisfy
every member of \(X\).  Then \(\mathcal B\) satisfies \((\mathrm{Ct})\),
\((\mathrm{Cl})\), \((\mathrm{Dj})\) and the parameter guard, and \(\sigma\)
respects \(\varnothing\).  By \cref{thm:p02-w02b-p02-033-restriction}, clause~1,
\(\mathcal B^{\downarrow}\) is strict and
\(\theta\defeq\theta_{\varnothing,\sigma}\) is a total assignment into it.
For each \(\gamma\in\Gamma\), \(\mathcal B,\sigma\entails
\operatorname{tr}^{\mathsf n}(\gamma)\), so clause~3 of that theorem gives
\(\mathcal B^{\downarrow}\entails\gamma[\theta]\), that is,
\(\mathcal B^{\downarrow},\theta\entails\gamma\).  The hypothesis gives
\(\mathcal B^{\downarrow},\theta\entails\varphi\), and clause~3 again gives
\(\mathcal B,\sigma\entails\operatorname{tr}^{\mathsf n}(\varphi)\).  So
\(X\entails_{\mathrm{FOL}}\operatorname{tr}^{\mathsf n}(\varphi)\).  This
step used neither \((\mathrm{Ne})\) nor \(\mathsf{Cov}\).

\emph{From one-sorted to many-sorted.}  Suppose
\(\Gamma\not\entails_{\mathrm{MS}}\varphi\): some strict \(\mathcal M\) and
total \(g\) satisfy every member of \(\Gamma\) and not \(\varphi\).  Let
\(\mathcal B\defeq\mathcal M^{\circ}_{g}\) if \(X\) contains \(\mathsf{Cov}\),
and \(\mathcal B\defeq\mathcal M^{\star}_{g}\) otherwise.  Fix
\(\mathsf s_0\in\mathsf{Sort}\) and let \(\sigma\) send every one-sorted
variable to \((\mathsf s_0,g(x^{\mathsf s_0}))\), an element of the domain of
either structure.  By \cref{thm:p02-w02b-p02-033-tagging}, clause~1,
\(\mathcal B\) satisfies \(\mathsf{Ax}_{\mathrm{sort}}\), hence
\(\mathsf{Ax}^{-}\), and the parameter guard, and
\(\mathcal M^{\circ}_{g}\) satisfies \(\mathsf{Cov}\) when \(\mathsf{Sort}\)
is finite.  By clause~3 of that theorem,
\(\mathcal B,\sigma\entails\operatorname{tr}^{\mathsf n}(\gamma)\) for every
\(\gamma\in\Gamma\) and
\(\mathcal B,\sigma\not\entails\operatorname{tr}^{\mathsf n}(\varphi)\).  So
\(\mathcal B,\sigma\) satisfies every member of \(X\) and refutes
\(\operatorname{tr}^{\mathsf n}(\varphi)\), and
\(X\not\entails_{\mathrm{FOL}}\operatorname{tr}^{\mathsf n}(\varphi)\).

The two steps prove the equivalence for each \(X\), hence clauses 1--3.  The
premise sets of clause~1 contain no coverage sentence, and when no coverage
sentence is present the countermodel \(\mathcal M^{\star}_{g}\) contains
\(\star\), which lies in no \(Q_{\mathsf s}\).  Every object used is defined
explicitly from \(\mathcal M\), \(g\), \(\mathcal B\), \(\sigma\) and the one
fixed sort \(\mathsf s_0\), so no step uses the Axiom of Choice.  No step
uses a bound on \(\abs{\mathsf{Sort}}\), \(\abs{\mathcal F}\) or
\(\abs{\mathcal R}\): the sort families, the parameter guard and
\(\operatorname{tr}^{\mathsf n}[\Gamma]\) are sets of sentences of
\(\Lang_{\mathrm{MS}}^{\mathsf n}\) of whatever size, and each translated
formula uses finitely many variables of \(\mathsf{Var}\)
(\cref{thm:p02-w02b-p02-033-sentences}).
\end{proof}
\end{annexitem}

\begin{annexitem}{annex:p02-w02b-p02-033-variable-form}{thm:p02-w02b-p02-033-variable-form}{The variable form for every injective renaming}
\begin{proof}
\emph{Bridge claim.}  Let \(\mathcal A\) be an
\(\Lang_{\mathrm{MS}}^{\flat}\)-structure, \(\mathcal A'\) an
\(\Lang_{\mathrm{MS}}^{\mathsf n}\)-expansion of it, \(\eta\) an environment,
\(\sigma\) an assignment into \(\mathcal A\) and \(\rho\) one into
\(\mathcal A'\), such that for every \(v\in\mathrm{Var}_{\mathrm{MS}}\)
\[
 \sigma(\iota(v))=\rho(\eta(v))\ \text{if }v\in\dom\eta,\qquad
 \sigma(\iota(v))=\mathsf n_v^{\mathcal A'}\ \text{if }v\notin\dom\eta.
 \tag{\(\dagger\)}
\]
Then \(\llbracket t^{\iota}\rrbracket^{\mathcal A}_{\sigma}=
\llbracket t^{\eta}\rrbracket^{\mathcal A'}_{\rho}\) for every term \(t\), and
\(\mathcal A,\sigma\entails\operatorname{tr}_{\iota}(\varphi)\) iff
\(\mathcal A',\rho\entails\operatorname{tr}^{\mathsf n}_{\eta}(\varphi)\) for
every well-sorted \(\varphi\).

For terms, by induction on \(t\): a variable \(v\) has value
\(\sigma(\iota(v))\) on the left and \(\rho(\eta(v))\) or
\(\mathsf n_v^{\mathcal A'}\) on the right, equal by \((\dagger)\); a
constant and a function symbol have the same interpretation in
\(\mathcal A\) and its expansion \(\mathcal A'\), so the remaining cases
follow from the induction hypothesis.  For formulas, by induction on
\(\varphi\), simultaneously for all \(\eta,\sigma,\rho\) satisfying
\((\dagger)\): relation atoms and equations follow from the term case,
because \(P^{\mathcal A}=P^{\mathcal A'}\); \(\bottom\) is false on both
sides; and each connective is the same truth function on both sides.  Let
\(\varphi=\forall v\,\chi\), \(\mathsf s\defeq\operatorname{srt}(v)\),
\(k\defeq k(\eta)\), \(\eta'\defeq\eta\langle v\rangle\), and note
\(\dom\mathcal A=\dom\mathcal A'\) and
\(Q_{\mathsf s}^{\mathcal A}=Q_{\mathsf s}^{\mathcal A'}\).  The left side
holds iff \(\mathcal A,\sigma_d\entails\operatorname{tr}_{\iota}(\chi)\) for
every \(d\in Q_{\mathsf s}^{\mathcal A}\), where
\(\sigma_d\defeq\sigma[\iota(v)\mapsto d]\); the right side holds iff
\(\mathcal A',\rho_d\entails\operatorname{tr}^{\mathsf n}_{\eta'}(\chi)\) for
every such \(d\), where \(\rho_d\defeq\rho[x_k\mapsto d]\).  Fix \(d\).  The
triple \(\eta',\sigma_d,\rho_d\) satisfies \((\dagger)\): at \(v\),
\(\sigma_d(\iota(v))=d=\rho_d(x_k)=\rho_d(\eta'(v))\); at \(u\ne v\),
injectivity of \(\iota\) gives \(\iota(u)\ne\iota(v)\), so
\(\sigma_d(\iota(u))=\sigma(\iota(u))\); if \(u\in\dom\eta'\), then
\(u\in\dom\eta\) and \(\eta'(u)=\eta(u)\ne x_k\), so
\(\rho_d(\eta'(u))=\rho(\eta(u))=\sigma(\iota(u))\); if
\(u\notin\dom\eta'\), then \(u\notin\dom\eta\) and
\(\sigma(\iota(u))=\mathsf n_u^{\mathcal A'}\).  The induction hypothesis
gives \(\mathcal A,\sigma_d\entails\operatorname{tr}_{\iota}(\chi)\) iff
\(\mathcal A',\rho_d\entails\operatorname{tr}^{\mathsf n}_{\eta'}(\chi)\), so
the two sides agree.  For \(\varphi=\exists v\,\chi\) the two sides say that
the same equivalent statements hold for some \(d\in Q_{\mathsf s}^{\mathcal A}\)
(by the clauses for \(\exists\) and \(\land\)), and the verification of
\((\dagger)\) for each \(d\) is the one just written, which does not depend on
the quantifier.

\emph{Clause 1.}  Take \(\eta=\varnothing\) and
\(\mathcal A'=\mathcal A_{\sigma}^{\iota}\).  Then \((\dagger)\) holds for
every \(\rho\), because \(\mathsf n_v^{\mathcal A_{\sigma}^{\iota}}=
\sigma(\iota(v))\).  The bridge claim gives the first equivalence for every
\(\rho\).  For the guards, \(\mathcal A,\sigma\entails
Q_{\operatorname{srt}(v)}(\iota(v))\) iff \(\sigma(\iota(v))\in
Q_{\operatorname{srt}(v)}^{\mathcal A}\) iff \(\mathsf n_v^{\mathcal A_{\sigma}^{\iota}}
\in Q_{\operatorname{srt}(v)}^{\mathcal A_{\sigma}^{\iota}}\) iff
\(\mathcal A_{\sigma}^{\iota}\entails Q_{\operatorname{srt}(v)}(\mathsf n_v)\).

\emph{Clause 2.}  Let \(Y\) be \(\mathsf{Ax}_{\mathrm{sort}}\), or
\(\mathsf{Ax}_{\mathrm{sort}}\cup\{\mathsf{Cov}\}\) when \(\mathsf{Sort}\) is
finite; \(Y\) is a set of \(\Lang_{\mathrm{MS}}^{\flat}\)-sentences.  By
\cref{thm:p02-w02b-p02-033-reduction}, clauses 1 and~2, it suffices to show
\[
 Y\cup\mathsf{VarGuard}_{\iota}\cup\operatorname{tr}_{\iota}[\Gamma]
 \entails_{\mathrm{FOL}}\operatorname{tr}_{\iota}(\varphi)
 \quad\text{iff}\quad
 Y\cup\mathsf{ParGuard}\cup\operatorname{tr}^{\mathsf n}[\Gamma]
 \entails_{\mathrm{FOL}}\operatorname{tr}^{\mathsf n}(\varphi).
\]
Right to left: let \(\mathcal A,\sigma\) satisfy the left premises.  By
\cref{lem:reduct-invariance}, \(\mathcal A_{\sigma}^{\iota}\entails Y\); by
clause~1 it satisfies the parameter guard and, at the assignment \(\sigma\),
every \(\operatorname{tr}^{\mathsf n}(\gamma)\).  So
\(\mathcal A_{\sigma}^{\iota},\sigma\entails
\operatorname{tr}^{\mathsf n}(\varphi)\), and clause~1 gives
\(\mathcal A,\sigma\entails\operatorname{tr}_{\iota}(\varphi)\).  Left to
right: let \(\mathcal B,\rho\) satisfy the right premises, and let
\(\mathcal A\) be the \(\Lang_{\mathrm{MS}}^{\flat}\)-reduct of
\(\mathcal B\).  Define \(\sigma(\iota(v))\defeq\mathsf n_v^{\mathcal B}\) for
\(v\in\mathrm{Var}_{\mathrm{MS}}\), well defined since \(\iota\) is injective,
and \(\sigma(x)\defeq\rho(x)\) for \(x\notin\operatorname{ran}\iota\).  Then
\(\mathcal A_{\sigma}^{\iota}=\mathcal B\): they share the reduct
\(\mathcal A\) and \(\mathsf n_v^{\mathcal A_{\sigma}^{\iota}}=
\sigma(\iota(v))=\mathsf n_v^{\mathcal B}\).  By \cref{lem:reduct-invariance},
\(\mathcal A\entails Y\); by clause~1, \(\sigma\) satisfies
\(\mathsf{VarGuard}_{\iota}\) and every \(\operatorname{tr}_{\iota}(\gamma)\).
So \(\mathcal A,\sigma\entails\operatorname{tr}_{\iota}(\varphi)\), and
clause~1 gives \(\mathcal B,\rho\entails\operatorname{tr}^{\mathsf n}(\varphi)\).

\emph{Clause 3, restriction.}  Let \(\mathcal B\defeq\mathcal A_{\sigma}^{\iota}\).
By \cref{lem:reduct-invariance} and clause~1, \(\mathcal B\) satisfies
\(\mathsf{Ax}_{\mathrm{sort}}\) and the parameter guard.
\(\mathcal B^{\downarrow}=\mathcal A^{\downarrow}\), since the restriction
reads only the \(\Lang_{\mathrm{MS}}^{\flat}\)-reduct.  For any assignment
\(\rho\), \(\theta_{\varnothing,\rho}(v)=\mathsf n_v^{\mathcal B}=
\sigma(\iota(v))\), so \(\theta_{\varnothing,\rho}=\sigma\circ\iota\).
\Cref{thm:p02-w02b-p02-033-restriction} gives that \(\mathcal A^{\downarrow}\)
is strict (nonempty carriers by its clause~1, disjointness by
\((\mathrm{Dj})\)), that \(\sigma\circ\iota\) is total into it, and that
\(\mathcal B,\rho\entails\operatorname{tr}^{\mathsf n}(\varphi)\) iff
\(\mathcal A^{\downarrow}\entails\varphi[\sigma\circ\iota]\).  Clause~1
replaces the left side by \(\mathcal A,\sigma\entails
\operatorname{tr}_{\iota}(\varphi)\).

\emph{Clause 3, tagging.}  \(g^{\sharp}\) is well defined because \(\iota\) is
injective.  The \(\Lang_{\mathrm{MS}}^{\flat}\)-reduct of
\(\mathcal M^{\star}_{g}\) is \(\mathcal M^{\sharp}\)
(\cref{def:p02-w02b-p02-033-tagged}), and
\(\mathsf n_v^{\mathcal M^{\star}_{g}}=(\operatorname{srt}(v),g(v))=
g^{\sharp}(\iota(v))\), so \((\mathcal M^{\sharp})^{\iota}_{g^{\sharp}}=
\mathcal M^{\star}_{g}\).  By \cref{thm:p02-w02b-p02-033-tagging}, clause~1,
and \cref{lem:reduct-invariance}, \(\mathcal M^{\sharp}\entails
\mathsf{Ax}_{\mathrm{sort}}\); by the same clause and clause~1 here,
\(g^{\sharp}\) satisfies \(\mathsf{VarGuard}_{\iota}\).  By clause~3 of the
tagging theorem and clause~1 here, \(\mathcal M,g\entails\varphi\) iff
\(\mathcal M^{\star}_{g},\rho\entails\operatorname{tr}^{\mathsf n}(\varphi)\)
iff \(\mathcal M^{\sharp},g^{\sharp}\entails\operatorname{tr}_{\iota}(\varphi)\).

\emph{The final sentence.}  For the fixed injection of
\cref{def:p02-many-sorted-signature}, \(\operatorname{tr}_{\iota}=
\operatorname{tr}_{\mathrm{MS}}\), \(\mathsf{VarGuard}_{\iota}=
\mathsf{VarGuard}_{\mathrm{MS}}\), \(\sigma\circ\iota\) is
\(\sigma{\restriction}\mathrm{Var}_{\mathrm{MS}}\) under the identification of
each sorted variable with its image, and \(g^{\sharp}\) is the assignment
\(s^{\sharp}\) of \cref{thm:p02-w02-p02-038-strict-correspondence}.  Clause~2
is then the statement of \cref{thm:p02-many-sorted-translation} and, with
\(\mathsf{Cov}\), of \cref{thm:p02-w02-p02-038-coverage-reduction}; clause~3
is clauses 1 and~2 of \cref{thm:p02-w02-p02-038-strict-correspondence}, whose
clause~3 is clause~2 here.  No step above refers to the cardinality of
\(\mathcal F\) or \(\mathcal R\).

\emph{Existence of injective renamings.}  If \(e:\mathsf{Sort}\to\N\) is
injective and each \(V_{\mathsf s}=\{x^{\mathsf s}_m:m\in\N\}\) without
repetition, put \(\iota(x^{\mathsf s}_m)\defeq x_{2^{e(\mathsf s)}(2m+1)-1}\).
If \(2^{a}(2m+1)=2^{b}(2m'+1)\), the largest powers of \(2\) dividing the two
sides agree, so \(a=b\), and then \(m=m'\).  So equal images force
\(e(\mathsf s)=e(\mathsf t)\), hence \(\mathsf s=\mathsf t\), and \(m=m'\),
hence the same variable.
\end{proof}
\end{annexitem}

\begin{annexitem}{annex:p02-w02b-p02-033-no-renaming}{thm:p02-w02b-p02-033-no-renaming}{No exact variable form with uncountably many sorts}
\begin{proof}
\emph{Clause 1.}  Suppose \(\iota\) is injective.  The map
\(\mathsf s\mapsto x^{\mathsf s}\) is injective, since distinct sorts have
disjoint stocks, and \(x_n\mapsto n\) is injective on \(\mathsf{Var}\).  The
composite \(\mathsf s\mapsto\iota(x^{\mathsf s})\mapsto\) its index is an
injection of \(\mathsf{Sort}\) into \(\N\), so \(\mathsf{Sort}\) is
countable.  Contrapositively, if \(\mathsf{Sort}\) is uncountable no
renaming is injective.

\emph{A strict structure.}  Let \(\mathcal M_0\) have carriers
\(M_{\mathsf s}\defeq\{\mathsf s\}\times\{0,1\}\), interpret
\(c:\longrightarrow\mathsf s\) as \((\mathsf s,0)\), let every function
symbol of codomain \(\mathsf s\) take the constant value \((\mathsf s,0)\),
and let every relation symbol be empty.  The carriers are nonempty and
pairwise disjoint, as their elements carry their sort as first coordinate, so
\(\mathcal M_0\) is strict.

\emph{Clause 2.}  Let \(u\ne w\) be sorted variables with
\(\iota(u)=\iota(w)\).

\emph{Case 1: \(\operatorname{srt}(u)=\operatorname{srt}(w)=\mathsf s\).}
Then \(\varphi\defeq(u=w)\) is well sorted.  Let \(g(u)\defeq(\mathsf s,0)\),
\(g(w)\defeq(\mathsf s,1)\) and \(g(v)\defeq(\operatorname{srt}(v),0)\) for
every other \(v\).  Then \(\mathcal M_0,g\not\entails u=w\), so
\(\varnothing\not\entails_{\mathrm{MS}}\varphi\).  But
\(\operatorname{tr}_{\iota}(\varphi)\) is \(\iota(u)=\iota(u)\), an equation
between one variable and itself, which every one-sorted structure satisfies
under every assignment.  So
\(\mathsf{Ax}_{\mathrm{sort}}\cup\mathsf{VarGuard}_{\iota}\entails_{\mathrm{FOL}}
\operatorname{tr}_{\iota}(\varphi)\).

\emph{Case 2: \(\operatorname{srt}(u)=\mathsf s\ne\mathsf t=
\operatorname{srt}(w)\).}  Let \(\varphi\defeq\bottom\).  With
\(g(v)\defeq(\operatorname{srt}(v),0)\) for every \(v\), \(\mathcal M_0,g\)
refutes \(\bottom\), so \(\varnothing\not\entails_{\mathrm{MS}}\bottom\).  If a
one-sorted \(\mathcal A,\sigma\) satisfied
\(\mathsf{Ax}_{\mathrm{sort}}\cup\mathsf{VarGuard}_{\iota}\), then
\(d\defeq\sigma(\iota(u))=\sigma(\iota(w))\) would lie in
\(Q_{\mathsf s}^{\mathcal A}\) and in \(Q_{\mathsf t}^{\mathcal A}\),
contradicting the \((\mathrm{Dj})\) sentence
\(\forall x\,\neg(Q_{\mathsf s}(x)\land Q_{\mathsf t}(x))\).  So the premise
set has no model, and it entails \(\operatorname{tr}_{\iota}(\bottom)=\bottom\)
because no structure and assignment satisfy the premises.

\emph{The final sentence.}  If \(\mathsf{Sort}\) is uncountable, every
renaming \(\iota\) is non-injective by clause~1, and clause~2 provides
\(\Gamma=\varnothing\) and \(\varphi\) on which the two sides of the
equivalence in clause~2 of \cref{thm:p02-w02b-p02-033-variable-form} differ.
\Cref{thm:p02-w02b-p02-033-reduction} holds for the same signature.
\end{proof}
\end{annexitem}
